\begin{figure*}[ht!]
    \vspace{-30mm}
    \includegraphics[width=\textwidth]{figures/architecture.pdf}
    \vspace{-30mm}
    \caption{
    Overview of our autoregressive diffusion transformer architecture. Tokens in the current frame only attend to tokens from current and previous frames, but not the future.
    \label{fig:ardit} 
    }
\end{figure*}


\section{Methods}
\label{sec:method}

Our approach introduces an autoregressive diffusion transformer that enables sequential video generation~(Sec.~\ref{sec:autoregressive_arch}). 
We show our training procedure in Fig.~\ref{fig:method}, which uses asymmetric distillation~(Sec.~\ref{sec:teacher}) and ODE initialization~(Sec.~\ref{sec:student_init}) to achieve high generation quality and stable convergence. We achieve efficient streaming inference through KV caching mechanisms~(Sec.~\ref{sec:kv_cache}).

\subsection{Autoregressive Architecture}
\label{sec:autoregressive_arch}
We begin by compressing the video into a latent space using a 3D VAE. The VAE encoder processes each chunk of video frames independently, compressing them into shorter chunks of latent frames. The decoder then reconstructs the original video frames from each latent chunk.
Our causal diffusion transformer operates in this latent space, generating latent frames sequentially. We design a block-wise causal attention mechanism inspired by prior works that combine autoregressive models with diffusion~\cite{lee2023diffusion,zhou2024transfusion,liu2024autoregressive,li2024autoregressive}.
Within each chunk, we apply bidirectional self-attention among latent frames to capture local temporal dependencies and maintain consistency.
To enforce causality, we apply causal attention across chunks. This prevents frames in the current chunk from attending to frames in future chunks.
A visual illustration of the architecture of our autoregressive diffusion transformer is provided in Fig.~\ref{fig:ardit}. 
Our design maintains the same latency as fully causal attention, as the VAE decoder still requires at least a block of latent frames to generate pixels. 
Formally, we define the attention mask $M$ as
\begin{equation} 
M_{i,j} = 
\begin{cases} 
1, & \text{if } \left\lfloor \dfrac{j}{k} \right\rfloor \leq \left\lfloor \dfrac{i}{k} \right\rfloor, \\
0, & \text{otherwise.} 
\end{cases} 
\end{equation}
Here, $i$ and $j$ index the frames in the sequence, $k$ is the chunk size, and $\left\lfloor \cdot \right\rfloor$ denotes the floor function. 

Our diffusion model $G_{\phi}$ extends the DiT architecture~\cite{peebles2023scalable} for autoregressive video generation. 
We introduce block-wise causal attention masks to the self-attention layers (illustrated in Fig.~\ref{fig:method}) while preserving the core structure, allowing us to leverage pretrained bidirectional weights for faster convergence.




\begin{figure*}
    \centering
    \includegraphics[width=\textwidth]{figures/methods_v2.pdf}
    \caption{
    Our method distills a many-step, bidirectional video diffusion model $s_\text{data}$ into a 4-step, causal generator $G_\phi$. The training process consists of two stages. (top) Student Initialization: we initialize the causal student by pretraining it on a small set of ODE solution pairs generated by the bidirectional teacher (Sec.~\ref{sec:student_init}). This step helps stabilize the subsequent distillation training.
    (bottom) Asymmetric Distillation: using the \textit{bidirectional} teacher, we train the \textit{causal} student generator through a distribution matching distillation loss~(Sec.~\ref{sec:teacher}).
    \label{fig:method} 
    }\vspace{-2em}
\end{figure*}



\subsection{Bidirectional \texorpdfstring{$\rightarrow$}{} Causal Generator Distillation}
\label{sec:teacher}

\begin{algorithm}[t]
\small 
\caption{Asymmetric Distillation with DMD \label{alg:training}}
\begin{algorithmic}[1]
\Require Few-step denoising timesteps $\mathcal{T} = \{0, t_1, t_2, \dots, t_Q \}$, video length $N$, chunk size $k$, pretrained bidirectional teacher model $s_{\text{data}}$, dataset $\mathcal{D}$.
\vspace{0.5em}
\State \textbf{Initialize} student model $G_\phi$ with ODE regression~(Sec.~\ref{sec:student_init})
\State \textbf{Initialize} generator output's score function $s_{\text{gen},\xi}$ with $s_{\text{data}}$

\While{training}
    \State \textbf{Sample} a video from the dataset and divide frames into $L = \lceil N / k \rceil$ chunks, $\{ x_0^{i} \}_{i=1}^L \sim \mathcal{D}$
    \State \textbf{Sample} per-chunk timesteps $\{ t^{i} \}_{i=1}^L\sim \text{Uniform}(\mathcal{T})$
    \State \textbf{Add noise:} ${ x_{t^i}^{i} = \alpha_{t^{i}} x_0^{i} + \sigma_{t^{i}} \epsilon^{i},\ \epsilon^{i} \sim \mathcal{N}(0, I)}$
    
    \State \textbf{Predict} clean frames with the student (block-wise causal mask): 
          $\hat{x}_0 = G_\phi\Bigl(\{ x_{t^i}^i \}, \{ t^i \}\Bigr)$\;

    
    \State \textbf{Sample} a single timestep $t\sim \text{Uniform}(0, T)$
    \State \textbf{Add noise} to predictions: $\hat{x}_t = \alpha_t \hat{x}_0 + \sigma_t \epsilon$, $\epsilon\sim \mathcal{N}(0, I)$
    
    \State \textbf{Update} student model $G_\phi$ using DMD loss $\mathcal{L}_{\text{DMD}}$ \Comment{Eq.~\ref{eq:dmd}}    
    
    \State \textbf{Train} generator's score function $s_{\text{gen},\xi}$:
    \State \quad \textbf{Sample new noise} $\epsilon'\sim \mathcal{N}(0, I)$
    \State \quad \textbf{Generate noisy} $\hat{x}_t$: $\hat{x}_t = \alpha_t \hat{x}_0 + \sigma_t \epsilon'$
    \State \quad Compute denoising loss and update $s_{\text{gen},\xi}$. \Comment{Eq.~\ref{eq:denoising_loss}}    
\EndWhile
\end{algorithmic}
\end{algorithm}


    
    
    
    


A straightforward approach to training a few-step causal generator would be through distillation from a causal teacher model. This involves adapting a pretrained bidirectional DiT model by incorporating causal attention mechanism described above and fine-tuning it using the denoising loss (Eq.~\ref{eq:denoising_loss}).
During training, the model takes as input a sequence of $N$ noisy video frames divided into $L$ chunks $\{x_t^{i}\}_{i=1}^{L}$, where $i \in \{1, 2, ..., L\}$ denotes the chunk index. Each chunk $x_t^{i}$ has its own noise time step $t^{i} \sim [0, 999]$, following Diffusion Forcing~\cite{chen2024diffusion}. 
During inference, the model denoises each chunk sequentially, conditioned on the previously generated clean chunks of frames.
While distilling this fine-tuned autoregressive diffusion teacher appears promising in theory, our initial experiments indicated that this naive approach yields suboptimal results. Since causal diffusion models typically underperform their bidirectional counterparts, training a student model from a weaker causal teacher inherently limits the student's capabilities. Moreover, issues such as error accumulation would propagate from teacher to student.
To overcome the limitations of a causal teacher, we propose an \textit{asymmetric distillation} approach: following state-of-the-art video models~\cite{videoworldsimulators2024,polyak2024movie}, we employ bidirectional attention in the teacher model while constraining the student model to causal attention~(Fig.~\ref{fig:method} bottom).
Algorithm~\ref{alg:training} details our training process.

















\subsection{Student Initialization}
\label{sec:student_init}


Directly training the causal student model using the DMD loss can be unstable due to architectural differences. 
To address this, we introduce an efficient initialization strategy to stabilize training (Fig.~\ref{fig:method} top).

We create a small dataset of ODE solution pairs generated by the bidirectional teacher model:

\begin{itemize} 
    \item Sample a sequence of noise inputs $\{x_T^{i}\}_{i=1}^{L}$ from the standard Gaussian distribution $\mathcal{L}(0, I)$. 
    \item Simulate the reverse diffusion process with an ordinary differential equation (ODE) solver~\cite{song2020denoising} using the pretrained bidirectional teacher model to obtain the corresponding ODE trajectory $\{x_t^{i}\}_{i=1}^{L}$, where $t$ spans $T$ to $0$, covering all inference timesteps.  
\end{itemize}


\noindent 
From the ODE trajectories, we select a subset of $t$ values that match those used in our student generator. 
The student model is then trained on this dataset with a regression loss:
\begin{equation} \mathcal{L}_{\text{init}} = \mathbb{E}_{x,t^i}  || G_{\phi}(\{x^i_{t^i}\}_{i=1}^N, \{t^i\}_{i=1}^N) - \{x^i_0\}_{i=1}^N ||^2\,,
\end{equation}
where the few-step generator $G_{\phi}$ is initialized from the teacher model.
Our ODE initialization is computationally efficient, requiring only a small number of training iterations on relatively few ODE solution pairs. 

\subsection{Efficient Inference with KV Caching}
\label{sec:kv_cache}

During inference, we generate video frames sequentially using our autoregressive diffusion transformer with KV caching for efficient computation~\cite{brown2020language}. 
We show the detailed inference procedure in Algorithm~\ref{alg:inference}. 
Notably, because we employ KV caching, block-wise causal attention is no longer needed at inference time. 
This allows us to leverage a fast bidirectional attention implementation~\cite{dao2022flashattention}.



    


    


\begin{algorithm}[t]
\small 
\caption{\label{alg:inference}Inference Procedure with KV Caching} 
\begin{algorithmic}[1]
\Require Denoising timesteps $\{t_0=0, t_1, \dots, t_Q \}$, video length $N$, chunk size $k$, few-step autoregressive video generator $G_{\phi}$, 

\State Divide frames into $L = \lceil N / k \rceil$ chunks
\State \textbf{Initialize} KV cache $\text{C} \gets \emptyset$

\For{$i = 1$ to $L$}
    \State \textbf{Initialize current chunk:} $x_{t_Q}^{i} \sim \mathcal{N}(0, I)$
    
    \State \textbf{Iterative denoising over timesteps:}
    \For{$j = Q$ to $1$}
        \State \textbf{Generate output:} $\hat{x}_{t_j}^{i} = G_{\phi}(x_{t_j}^{i}, t_j)$ using cache $\text{C}$
        \State \textbf{Update chunk:} $x_{t_{j-1}}^{i} = \alpha_{t_{j-1}} \hat{x}_{t_j}^{i} + \sigma_{t_{j-1}} \epsilon'$ %
    \EndFor
    
    \State \textbf{Update KV cache:}
    \State Compute KV pairs with a forward pass $G_{\phi}(x_{0}^{i}, 0)$
    \State Append new KV pairs to cache $\text{C}$
\EndFor

\State \textbf{Return} $\{ x_{0}^{i} \}_{i=1}^{L}$
\end{algorithmic}
\end{algorithm}





    


    



